\documentclass[10pt]{article}
\usepackage{icmlko}

\icmlmeta{IP 파운데이션 월드모델}{503}{대리변수가 아니었다}{도메인 이름표가 층 ③ 을 대신할 수 있다는 결론은 464행 표의 성질이었다. 자를 고정한 채 표만 키우니 부호가 뒤집힌다}

\begin{document}

\twocolumn[
\icmlheader

\begin{abstract}
\textbf{이 논문은 두 가지를 한다.} 첫째, 앞 스텝(502)이 \textbf{자 하나로} 찍은 판정을
\textbf{네 자로} 다시 낸다. 502 는 팔 B $=$ 2{,}050 에서 $\Delta\rho($층 ③$)=+$0.058238 이
문턱을 넘는 것만 보고 「든다」를 실었다. 부 표적의 부호($\text{㉡}$)도, 순열 영분포($\text{㉢}$)도,
\emph{같은 사이클이 스스로 만든} 「더 싼 대체물」 조건($\text{㉣}$)도 대지 않았다. 넷을 다 대니
$\text{㉠}$ $+$0.058238 $>$ $T$ $=$ 0.020060, $\text{㉡}$ $+$0.005669(부호 유지),
$\text{㉢}$ 관측 $+$0.057857 $>$ 상위 5\,\% 지점 $-$0.000493, $\text{㉣}$ 도메인 지시자
10열이 혼자 들지만($+$0.096085) \textbf{그 위에서 층 ③ 도 산다}($+$0.007112 $>$ $T$ $=$
0.006098) --- \textbf{채택된다}. 둘째, \textbf{그 마지막 칸이 표 크기에 매달려 있음을 보인다.}
같은 코드 $\cdot$ 같은 재표집 수 $\cdot$ 같은 씨앗으로 464행 표를 다시 재면 같은 칸이
\textbf{$-$0.007687 로 음수}이고 판정은 「뺀다」다. \textbf{자는 고정됐고 움직인 것은 표뿐이다.}
그러므로 「층 ③ 의 정보는 도메인 이름표로 \emph{대체 가능}하다」(스텝 501 의 다섯째 주장)는
\textbf{464행 표의 성질}이었다. 도메인 구성 이동을 통제해도 부호는 양수로 남지만,
851행으로 줄이면 세 씨앗 중 둘이 「못 가른다」로 돌아간다 --- \textbf{이 결론은 2{,}050행을
필요로 한다}.
\end{abstract}
]

\section{자가 하나였다}

502 의 판정 경로는 \verb|verdict(dr, T)| 하나다. 그런데 같은 사이클이 저장소에
\verb|lab/adopt.py| 를 새로 넣고 그 문서에 「\textbf{부르는 쪽이 이 함수의 답을 키로 내야
한다}」고 적었다. 실제로 그 함수를 부르는 곳은 저장소 전체에서 \textbf{한 군데}였고, 그
자리는 \emph{앞 사이클을 채점하는} 자리였다. 자기 헤드라인에는 안 댔다.

\begin{center}\small
\begin{tabular}{llr}
\hline
자 & 무엇 & 값 \\
\hline
$\text{㉠}$ & 합산 $\Delta\rho > T$ & $+$0.058238 $/$ 0.020060 \\
$\text{㉡}$ & 부 표적 부호 유지 & $+$0.005669 \\
$\text{㉢}$ & 열 순열 영분포 밖 & $+$0.057857 $/$ $-$0.000493 \\
$\text{㉣}$ & 더 싼 대체물이 있나 & 없다 \\
\hline
\multicolumn{2}{l}{\textbf{채택}} & \textbf{참} \\
\hline
\end{tabular}
\end{center}

$\text{㉣}$ 을 \emph{안 쟀을 때} 같은 함수는 「모른다 --- 보류한다」를 낸다. 「모른다」는
「통과」가 아니다. \textbf{즉 502 의 헤드라인은 자기 저장소의 규칙으로는 그 시점에 보류
상태였다.} 이 논문은 그 보류를 측정으로 푼다.

\claimx{판정은 자의 개수에 딸린다. 같은 $\Delta\rho$ 라도 자 하나($\text{㉠}$)면 「든다」,
넷($\text{㉠㉡㉢㉣}$)이면 「채택」, $\text{㉣}$ 을 안 재면 「보류」다 --- 세 문장은 다르다.}
{$\text{㉡}$ 의 부호가 뒤집히거나 $\text{㉢}$ 이 영분포 안에 들어오면 채택이 죽는다.
실측은 각각 $+$0.005669 와 $+$0.057857 $>$ $-$0.000493 이다.}

\section{대체 가능성은 표 크기의 함수였다}

핵심 실험은 단순하다. \textbf{자를 완전히 고정하고 표만 바꾼다.} 같은 파일의 같은
함수, 재표집 4{,}000회, 씨앗 7, 군집 열 동일.

\begin{center}\small
\begin{tabular}{lrrl}
\hline
표 & 지시자 열 & $\Delta\rho$ & 판정 \\
\hline
464 & 9 & $-$0.007687 & 뺀다 \\
2{,}050 & 10 & $+$0.007112 & 든다 \\
\hline
\end{tabular}
\end{center}

464행 쪽 수는 스텝 501 이 적은 값($-$0.0075)과 소수 넷째 자리까지 맞고, 지시자 자신의
이득도 $+$0.087630 으로 그대로 재현된다. \textbf{재현이 되므로 「자가 달라졌다」는 설명은
죽는다.}

남는 설명은 둘이다 --- \textbf{표본 크기}와 \textbf{도메인 구성}. 표가 4.42배 크는 동안 한
도메인의 비중이 47.41\,\% 에서 55.02\,\% 로 올랐고, 층 ③ 은 \emph{도메인 배경 대비 편차}라
그 이동이 값에 섞인다. 그래서 2{,}050 을 464 의 도메인 비중으로 \textbf{851행 부분표집}해
다시 쟀다.

\begin{center}\small
\begin{tabular}{lrrl}
\hline
씨앗 & $\Delta\rho$ & $T$ & 판정 \\
\hline
11 & $+$0.00278 & 0.00377 & 못 가른다 \\
12 & $+$0.00978 & 0.00930 & 든다 \\
13 & $+$0.00203 & 0.00485 & 못 가른다 \\
\hline
\end{tabular}
\end{center}

\textbf{부호는 셋 다 양수다} --- 464 의 음수는 구성 이동의 산물이 아니다. \textbf{그러나
851행에서는 셋 중 둘이 「못 가른다」다.} 결론은 표 크기에 걸려 있고, 이 논문은 그것을
숨기지 않는다. 그리고 부분표집은 층 ③ 의 \emph{배경 풀}(3{,}896 개체)을 안 줄이므로
「표본이 커진 것」과 「배경이 좋아진 것」은 여전히 못 가른다.

\claimx{「층 ③ 의 정보는 도메인 이름표로 대체 가능하다」(스텝 501 의 다섯째 주장)를
철회한다. 그 문장은 464행 표에서만 참이고, 자를 고정한 채 2{,}050행으로 가면 부호가
뒤집힌다($-$0.007687 $\to$ $+$0.007112).}
{851행 부분표집에서 세 씨앗의 부호가 갈리거나 음수가 나오면 이 철회가 죽는다.
실측은 셋 다 양수이나 셋 중 둘이 문턱 아래다 --- \textbf{2{,}050행 밑으로는 못 선다}.}

\section{「군집 부트」라는 이름을 쓸 자격}

502 는 SE 를 「군집 부트」라 불렀다. 그 부트의 군집 열은 개체인데, 이 자료는
\textbf{개체당 정확히 한 행}이다.

\begin{center}\small
\begin{tabular}{lr}
\hline
팔 B 행 & 2{,}050 \\
서로 다른 개체 & 2{,}050 \\
서로 다른 도메인 & 10 \\
\hline
\end{tabular}
\end{center}

군집 수가 행 수와 같으면 그것은 \textbf{보통의 행 부트}다. 이름을 고친다. 그리고 같은
사이클이 「이득의 43.62\,\% 가 도메인 \emph{안}에 있다」를 보였으므로 행은 독립이 아니다.
도메인을 군집으로 놓으면

\[
T:\ 0.020060 \;\longrightarrow\; 0.053524 \quad (2.67\text{배})
\]

이고, \textbf{그 아래서도} $\Delta\rho = +$0.058238 이 선다(여유 1.088배). 부트 씨앗을
7 $\cdot$ 17 $\cdot$ 101 $\cdot$ 2026 으로 흔들어도 넷 다 「든다」이고 95\,\% 구간
$[+0.0290,\ +0.1416]$ 은 0 을 안 품는다. \textbf{다만 군집이 10개뿐이라 이 SE 는 아래로
치우친다 --- 이 $T$ 는 하한으로 읽어야 한다.}

\claimx{개체를 군집으로 쓴 부트는 군집이 없는 행 부트다(군집 수 $=$ 행 수 $=$ 2{,}050).
도메인을 군집으로 놓으면 문턱이 2.67배로 커지고, 커진 문턱 아래서도 판정은 산다.}
{도메인 군집 $T$ 가 0.058238 을 넘으면 「든다」가 죽는다. 실측 0.053524 로 여유는 8.8\,\%뿐이고,
군집 10개짜리 부트가 SE 를 아래로 치우치게 하므로 \textbf{이 여유는 상한이다}.}

\section{등록한 수를 쓰는 것}

502 의 사전등록은 재표집 \textbf{4{,}000}회를 등록하고 2{,}000회를 명시적으로 기각했다.
실제 코드는 400회였고 어디에도 신고가 없었다. 등록값으로 되돌려 보았다.

\[
\mathrm{SE}:\ 0.010137 \;\longrightarrow\; 0.010030 \quad (1.1\,\%)
\]

\textbf{수는 맞았다.} 그러나 「맞았으니 괜찮다」는 사후의 문장이고, 등록 시점에는 알 수
없었다. 그래서 이 사이클은 배선 검사를 하나 늘렸다 --- \textbf{러너가 사전등록 문서의
상수 블록을 직접 읽어 자기 모듈 상수와 기계로 맞대고, 하나라도 어긋나면 측정 전에
멈춘다}(16개 중 16개 일치).

\section{문턱은 무엇으로 주는가}

팔 A(물리 장)는 38행에 머물러 어떤 자로도 못 가른다($T = $ 0.2866). 앞 사이클은
「좌표를 118행에 붙이면 $T \approx 0.169$」라 적었고($T \propto 1/\sqrt{n}$), 비평은
「군집이 격자이고 격자는 23개뿐이니 $T \approx 0.28$」이라 적었다. \textbf{둘 다 안 재고
세운 식이다.} 38행 위에서 직접 흔들었다.

\begin{center}\small
\begin{tabular}{lrr}
\hline
흔든 것 & $\beta$ & $R^2$ \\
\hline
군집 수(행 고정) & $-$0.057 & 0.29 \\
행 수(격자 21 고정) & $+$0.706 & 0.75 \\
\hline
\end{tabular}
\end{center}

$\mathrm{SE} \propto x^{-\beta}$ 로 놓았다. \textbf{SE 는 군집 수에 사실상 반응하지 않고
행 수로 준다} --- 그것도 $1/\sqrt{n}$ 보다 빠르게. 실측 $\beta$ 로 「118행 $\cdot$ 23격자」의
문턱을 예측하면 \textbf{0.1294} 다.

\claimx{팔 A 의 문턱은 군집 수가 아니라 행 수로 준다($\beta_{\text{군집}} = -0.057$,
$\beta_{\text{행}} = +0.706$). 「118행 $\cdot$ 23격자」의 예측 문턱은 0.1294 이고, 이는
$1/\sqrt{n}$ 식의 0.1626 보다도 낮다.}
{38행에서 3.1배 밖으로 외삽한 값이고 행 곡선이 단조롭지 않다(0.2034 $\to$ 0.2207 $\to$
0.1540 $\to$ 0.1433). 실제로 118행을 만들었을 때 $T$ 가 0.25 를 넘으면 이 예측은 죽는다.}

\section{한 문장이 네 개의 수였다}

앞 사이클들이 「815 중 810」을 한 분모처럼 인용했다. 실측하면 네 개다.

\begin{center}\small
\begin{tabular}{lr}
\hline
무엇 & 값 \\
\hline
표적 함수가 낸 행 & 815 \\
그중 서로 다른 키 & 813 \\
그 주행의 성공 건수 & 810 \\
살아 있는 파일의 개체 수 & 813 \\
살아 있는 파일의 줄 수 & 815 \\
\hline
\end{tabular}
\end{center}

마지막 두 줄의 차 2 는 \textbf{같은 키가 두 번 실린 행}이다 --- 표적 함수가 키를 두 번
내보내서 두 번 받아 두 번 썼다. 「815 라는 개체」는 없다.

\claimx{하나의 인용문 「815 중 810」은 서로 다른 네 개의 분모를 겹쳐 놓은 것이다
(행 815 $\cdot$ 키 813 $\cdot$ 성공 810 $\cdot$ 파일 개체 813). 그리고 파일에는 중복 행이
2개 있다.}
{어느 한 쌍이라도 같은 수로 밝혀지면 이 주장이 약해진다. 실측은 넷이 서로 다르고
중복 키 둘이 이름으로 특정된다.}

\section{한계}

\textbf{첫째}, 「표본이 커진 것」과 「배경이 좋아진 것」을 못 가른다. 표를 키우면 층 ③ 의
배경 풀도 같이 커지고, 부분표집으로는 배경이 안 줄어든다. \textbf{둘째}, 도메인 군집
부트의 군집이 10개뿐이라 그 문턱은 하한이다. \textbf{셋째}, 팔 A 의 $\beta$ 는 38행에서
외삽한 값이고 씨앗 셋의 산포만 있다. \textbf{넷째}, 한 도메인(영화)의 406개 키가 위키
파이프라인에 원리상 안 들어간다 --- 구조가 있는 빈칸이고 이 사이클은 그 크기를 안 쟀다.
\textbf{다섯째}, 채택 규칙 자체에 구멍이 남아 있다. 「못 가른다」와 「안 든다」를 한
불리언으로 접어서, 도메인 군집으로 재면 두 하위 검사가 모두 「못 가른다」인데도
「통과」가 나온다. 이 사이클은 그 사실을 신고하고 고치지 않았다 --- 수리 상한을 측정
전에 6으로 선언했고 여섯을 이미 썼기 때문이다.

\end{document}
